\chapter{Introdu\c{c}\~ao}
\label{cap:introducao}

O contexto escolar produz diariamente um volume expressivo de dados acad\^emicos, administrativos e operacionais. Notas, faltas, registros financeiros, turmas, s\'eries e hist\'oricos disciplinares s\~ao continuamente armazenados pelos sistemas institucionais, mas nem sempre se convertem em informa\c{c}\~ao \'util para a tomada de decis\~ao pedag\'ogica. Na pr\'atica, a escola registra muito, mas ainda explora pouco esse conjunto de dados para antecipar dificuldades de aprendizagem e organizar interven\c{c}\~oes mais precoces. Este trabalho parte dessa lacuna.

Na rotina escolar, boa parte das interven\c{c}\~oes ainda acontece de forma reativa. Em geral, professores, coordena\c{c}\~ao e secretaria percebem o problema quando o aluno j\'a apresenta queda acentuada de desempenho, aus\^encias recorrentes ou risco concreto de reprova\c{c}\~ao. Nesse momento, a margem de a\c{c}\~ao \'e menor e o custo institucional tende a ser maior. Diante disso, este \sigla{TCC}{Trabalho de Conclus\~ao de Curso} investiga se os dados hist\'oricos j\'a dispon\'iveis na escola podem ser utilizados para antecipar sinais de risco pedag\'ogico com anteced\^encia suficiente para apoiar interven\c{c}\~oes reais.

O problema central deste trabalho pode ser sintetizado da seguinte forma: como integrar dados escolares hist\'oricos de notas, faltas, pagamentos e registros acad\^emicos em uma pipeline preditiva capaz de antecipar a pr\'oxima nota do aluno e sinalizar risco pedag\'ogico, produzindo relat\'orios \'uteis para a tomada de decis\~ao institucional? A partir dessa formula\c{c}\~ao, o trabalho foi conduzido pela hip\'otese de que uma arquitetura anal\'itica temporal, sustentada por dados tabulares e valida\c{c}\~ao cronol\'ogica, seria mais aderente ao problema real do que abordagens puramente relacionais ou experimentos isolados. Essa hip\'otese n\~ao foi apenas assumida de in\'icio: ela foi sendo confirmada ao longo da evolu\c{c}\~ao do projeto, especialmente quando a solu\c{c}\~ao deixou de priorizar grafos e passou a concentrar esfor\c{c}os em modelagem supervisionada sobre dados longitudinais.


\section{Objetivos}
\label{sec:objetivos}

O objetivo geral do trabalho \'e desenvolver uma pipeline preditiva demonstrada com dados escolares sintéticos para prever a pr\'oxima nota do aluno, identificar risco pedag\'ogico e gerar relat\'orios operacionais de apoio \`a decis\~ao para professores, coordena\c{c}\~ao e secretaria.

Como objetivos espec\'ificos, o trabalho busca:

\begin{enumerate}
    \item investigar a viabilidade de diferentes abordagens computacionais para o dom\'inio escolar;
    \item estruturar um fluxo de prepara\c{c}\~ao local, gera\c{c}\~ao de dados sintéticos e extra\c{c}\~ao compat\'ivel com o contrato p\'ublico;
    \item integrar dados de notas, faltas, pagamentos e contexto acad\^emico em uma base longitudinal;
    \item construir atributos temporais que representem a trajet\'oria recente do aluno;
    \item comparar modelos e baselines com valida\c{c}\~ao temporal;
    \item transformar as sa\'idas t\'ecnicas em relat\'orios compreens\'iveis para a escola.
\end{enumerate}

\section{Justificativa}
\label{sec:justificativa}

O trabalho se justifica em tr\^es planos complementares. No plano acad\^emico, dialoga com as \'areas de \sigla{EDM}{Educational Data Mining} e \sigla{LA}{Learning Analytics}, mostrando como um problema institucional real pode ser tratado por meio de integra\c{c}\~ao de dados, modelagem e gera\c{c}\~ao de artefatos operacionais. No plano tecnol\'ogico, a contribui\c{c}\~ao est\'a em sair de experimentos pontuais e consolidar uma pipeline reprodut\'ivel, com separa\c{c}\~ao entre prepara\c{c}\~ao local dos dados, modelagem e sa\'idas de uso institucional. No plano social e educacional, a relev\^ancia est\'a na possibilidade de apoiar interven\c{c}\~oes mais precoces, antes que o aluno atinja um quadro de maior gravidade.

Tamb\'em existe justificativa \'etica. Em educa\c{c}\~ao, sistemas preditivos n\~ao devem substituir o julgamento pedag\'ogico humano, mas sim apoiar a prioriza\c{c}\~ao e a leitura do contexto. Por isso, o sistema proposto foi concebido como ferramenta de apoio \`a decis\~ao e n\~ao como mecanismo autom\'atico de julgamento. Em outras palavras, o objetivo n\~ao \'e rotular o aluno de forma definitiva, mas oferecer sinais adicionais para que a escola consiga agir antes.

\section{Delimita\c{c}\~ao do tema}
\label{sec:delimitacao}

Este trabalho n\~ao trata de evas\~ao escolar em sentido amplo, nem pretende substituir os processos formais de avalia\c{c}\~ao da escola. O foco est\'a na previs\~ao da pr\'oxima nota e na identifica\c{c}\~ao de risco pedag\'ogico de curto prazo, com base em registros escolares representados por dados sintéticos. A pesquisa tamb\'em se limita ao contrato p\'ublico desses dados e \`as rotinas de prepara\c{c}\~ao e extra\c{c}\~ao que sustentam a pipeline publicada. Isso significa que as conclus\~oes do trabalho devem ser lidas dentro desse contexto espec\'ifico de dados, regras e disponibilidade institucional.

\section{Estrutura do trabalho}
\label{sec:estrutura}

Esta monografia foi organizada em seis cap\'itulos. O Cap\'itulo~\ref{cap:introducao} apresenta o problema, os objetivos e a justificativa. O Cap\'itulo~\ref{cap:referencial} discute a fundamenta\c{c}\~ao te\'orica e os trabalhos relacionados. O Cap\'itulo~\ref{cap:metodologia} descreve a natureza da pesquisa e a evolu\c{c}\~ao hist\'orica do projeto. O Cap\'itulo~\ref{cap:arquitetura} detalha a arquitetura proposta, a integra\c{c}\~ao dos dados e a constru\c{c}\~ao do dataset. O Cap\'itulo~\ref{cap:resultados} apresenta a modelagem, os relat\'orios e os resultados experimentais. Por fim, o Cap\'itulo~\ref{cap:conclusao} discute limita\c{c}\~oes, conclus\~oes e trabalhos futuros.
